\documentclass[12pt]{article}
\usepackage[utf8]{inputenc}
\usepackage[T1]{fontenc}
\usepackage[margin=1in]{geometry}
\usepackage{setspace}
\usepackage{xcolor}
\usepackage[colorlinks=true, urlcolor=blue!60!black, linkcolor=blue!60!black]{hyperref}
\setlength{\parindent}{2em}
\setlength{\parskip}{1.25em}
\renewcommand{\baselinestretch}{1.0}
\begin{document}
\begin{center}
{\LARGE \textit{Craig v. Boren}, 429 U.S. 190 (1976)\par}
\vspace{0.8em}
{\large Case Brief\par}
\end{center}
\vspace{1em}

\section{Memory Jogger}

In \href{https://www.courtlistener.com/opinion/109570/craig-v-boren/}{\textit{Craig v. Boren}}, the Supreme Court struck down an Oklahoma law barring the sale of 3.2\% beer to men under 21 while permitting sales to women at 18, formally establishing ``intermediate scrutiny'' for sex-based classifications under the Equal Protection Clause.

\section{Detailed Case Facts}

Two interacting provisions of Oklahoma's liquor code, Okla. Stat., Tit. 37, \S\S\ 241 and 245 (1958 and Supp. 1976), made it unlawful for any licensed vendor to sell, barter, or give ``nonintoxicating'' 3.2\% beer (defined as a beverage containing more than 0.5\% alcohol by volume and not more than 3.2\% alcohol by weight) to a ``minor,'' which \S\ 245 defined as a female under age 18 and a male under age 21 (\href{https://supreme.justia.com/cases/federal/us/429/190/}{429 U.S. at 191--192 \& n.1}). Notably, the statutory scheme regulated only the \textit{sale} of 3.2\% beer to 18--20-year-old men, not their \textit{possession} or \textit{consumption} once obtained---meaning an 18-year-old woman could lawfully buy 3.2\% beer and hand it to an 18--20-year-old male companion.

Before 1972, Oklahoma had set the general age of civil majority at 18 for females and 21 for males (Okla. Stat., Tit. 15, \S\ 13), while holding females criminally responsible as adults at 18 and males at 16 (Okla. Stat., Tit. 10, \S\ 1101(a)). After the Tenth Circuit invalidated the criminal-responsibility age differential under \textit{Reed v. Reed}, 404 U.S. 71 (1971), in \textit{Lamb v. Brown}, 456 F.2d 18 (10th Cir. 1972), the Oklahoma Legislature standardized both criminal responsibility and general civil majority at age 18 for both sexes in 1972---yet simultaneously carved out \S\S\ 241 and 245 to preserve the 21/18 gender gap for 3.2\% beer (\href{https://supreme.justia.com/cases/federal/us/429/190/}{429 U.S. at 197}).

On December 20, 1972, Curtis Craig---a male Oklahoma State University undergraduate then between 18 and 21 years old (who replaced original plaintiff Mark Walker after Walker turned 21)---and Carolyn Whitener, a licensed vendor of 3.2\% beer who co-owned the ``Honk-N-Holler'' convenience store in Stillwater, Oklahoma, sued Oklahoma Governor David Boren (named in his official capacity) and other state officials seeking declaratory and injunctive relief under the Equal Protection Clause of the Fourteenth Amendment (\href{https://ojs.library.okstate.edu/osu/index.php/OKPolitics/article/view/930/832}{Oklahoma Politics, 47--55}; \href{https://supreme.justia.com/cases/federal/us/429/190/}{429 U.S. at 192}). Oklahoma defended the law on two main grounds: (1) statistical surveys purportedly linking young males to higher rates of drunk driving and traffic casualties, and (2) broad state regulatory authority over alcoholic beverages under Section 2 of the Twenty-first Amendment.

Historically and politically, \textit{Craig} arrived at a decisive moment in 1970s sex-equality litigation. In \textit{Reed v. Reed} (1971), the Court had invalidated a sex-based probate preference using nominal rational-basis language. Two years later, a four-Justice plurality in \textit{Frontiero v. Richardson}, 411 U.S. 677 (1973), urged treating sex as a suspect classification subject to strict scrutiny, falling one vote short while the Equal Rights Amendment (ERA) remained pending before state legislatures. Ruth Bader Ginsburg, then leading the ACLU Women's Rights Project, advised plaintiffs' counsel, submitted an \textit{amicus} brief, and sat at counsel table during oral argument, strategically using a male plaintiff's challenge to an archaic sex stereotype to cement heightened scrutiny for all gender classifications (\href{https://www.oyez.org/cases/1976/75-628}{Oyez}).

\section{Procedural History}

\begin{itemize}
  \item \textbf{United States District Court for the Western District of Oklahoma:} Because plaintiffs sought to enjoin enforcement of a state statute on federal constitutional grounds, a three-judge District Court was convened under 28 U.S.C. \S\ 2281. In \textit{Walker v. Hall}, 399 F. Supp. 1304 (W.D. Okla. 1975), the court recognized \textit{Reed v. Reed} as controlling but upheld Oklahoma's statutory differential and dismissed the complaint, concluding that the State's statistical evidence on youthful male drunk-driving arrests and traffic injuries demonstrated a rational and substantial relationship to highway safety (\href{https://supreme.justia.com/cases/federal/us/429/190/}{429 U.S. at 192, 199--201}).
  \item \textbf{Basis for Supreme Court Review:} Appellants took a direct appeal to the U.S. Supreme Court, which noted probable jurisdiction sub nom. \textit{Craig v. Boren}, 423 U.S. 1047 (1976). After probable jurisdiction was noted, Curtis Craig turned 21, rendering the controversy moot as to him because no class certification had been sought and only declaratory and injunctive relief were requested (\href{https://supreme.justia.com/cases/federal/us/429/190/}{429 U.S. at 192 \& n.2}).
  \item \textbf{Disposition:} Reversed (7--2) on December 20, 1976. The Court held that vendor Whitener had third-party standing and invalidated Okla. Stat., Tit. 37, \S\ 245 as an unconstitutional denial of equal protection to males aged 18--20 (\href{https://tile.loc.gov/storage-services/service/ll/usrep/usrep429/usrep429190/usrep429190.pdf}{429 U.S. at 210}).
\end{itemize}

\section{Judicial Votes}

\begin{itemize}
  \item \textbf{Majority (7--2 in judgment; 6--2 in opinion except Part II-D):} Justice William J. Brennan, Jr. delivered the opinion of the Court, joined in full by Justices Byron R. White, Thurgood Marshall, Lewis F. Powell, Jr., and John Paul Stevens, and joined in all except Part II-D (the Twenty-first Amendment analysis) by Justice Harry A. Blackmun (\href{https://supreme.justia.com/cases/federal/us/429/190/}{429 U.S. at 191}).
  \item \textbf{Concurrences:}
    \begin{itemize}
      \item Justice Lewis F. Powell, Jr. filed a concurring opinion (joining the Court's opinion).
      \item Justice John Paul Stevens filed a concurring opinion (joining the Court's opinion).
      \item Justice Harry A. Blackmun filed a brief statement concurring in part (joining all but Part II-D of the Court's opinion).
      \item Justice Potter Stewart filed an opinion concurring in the judgment.
    \end{itemize}
  \item \textbf{Dissents:}
    \begin{itemize}
      \item Chief Justice Warren E. Burger filed a dissenting opinion.
      \item Justice William H. Rehnquist filed a dissenting opinion.
    \end{itemize}
\end{itemize}

\section{Holding}

\textbf{Judgment:} Reversed the judgment of the three-judge District Court, invalidating the gender-based age differential in Okla. Stat., Tit. 37, \S\S\ 241 and 245 (\href{https://supreme.justia.com/cases/federal/us/429/190/}{429 U.S. at 210}).

\textbf{Rules Announced:}
\begin{enumerate}
  \item \textbf{Mootness and Vendor \textit{Jus Tertii} Standing:} While a male challenger's request for prospective relief becomes moot when he turns 21, a licensed alcohol vendor directly regulated by the statute suffers Article III injury-in-fact (forced to choose between lost sales and regulatory sanctions) and has third-party (\textit{jus tertii}) standing to assert the equal protection rights of her excluded male customers (\href{https://supreme.justia.com/cases/federal/us/429/190/}{429 U.S. at 192--197}).
  \item \textbf{Intermediate Scrutiny for Sex Classifications:} Under the Equal Protection Clause, ``classifications by gender must serve important governmental objectives and must be substantially related to achievement of those objectives'' (\href{https://supreme.justia.com/cases/federal/us/429/190/}{429 U.S. at 197}).
  \item \textbf{Insufficiency of Weak Statistical Proxies:} Although traffic safety is an important governmental objective, arrest statistics showing that 2\% of 18--20-year-old males versus 0.18\% of females were arrested for driving under the influence provide far too tenuous a fit to justify using sex as a proxy for drinking and driving, particularly when the law prohibits only purchase and not consumption (\href{https://supreme.justia.com/cases/federal/us/429/190/}{429 U.S. at 199--204}).
  \item \textbf{Twenty-first Amendment Does Not Dilute Equal Protection:} Section 2 of the Twenty-first Amendment primarily qualifies the Dormant Commerce Clause and neither overrides nor relaxes Fourteenth Amendment Equal Protection scrutiny of invidious gender classifications (\href{https://supreme.justia.com/cases/federal/us/429/190/}{429 U.S. at 204--210}). Insofar as \textit{Goesaert v. Cleary}, 335 U.S. 464 (1948), suggested otherwise, it was disapproved (\href{https://supreme.justia.com/cases/federal/us/429/190/}{429 U.S. at 210 n.23}).
\end{enumerate}

\textbf{Holding vs. Dicta:} The Court's rulings on vendor third-party standing, the intermediate scrutiny formula, the invalidity of Oklahoma's 21/18 3.2\% beer statute, and the non-insulation of sex discrimination under the Twenty-first Amendment are essential holdings. By contrast, Justice Brennan's footnote 7 reservation of whether courts must accept post hoc litigation rationales offered by state attorneys, footnote 22's discussion of hypothetical racial and ethnic alcohol-consumption statistics, and footnote 24's observation that Oklahoma remained free to enact a uniform, gender-neutral drinking age are dicta.

\section{Analysis of Opinions}

\subsection{Majority Opinion (Justice Brennan)}

Justice Brennan's opinion proceeds in three major parts: standing, Equal Protection scrutiny, and the Twenty-first Amendment (\href{https://supreme.justia.com/cases/federal/us/429/190/}{429 U.S. at 192--210}). On standing (Part I), Brennan noted that Craig's claim was moot after his 21st birthday, but held that Whitener could press the equal protection claim both because Oklahoma had failed to object below (forfeiting prudential \textit{jus tertii} objections) and independently under \textit{Eisenstadt v. Baird}, 405 U.S. 438 (1972). Because \S\S\ 241 and 245 directly targeted vendors---forcing Whitener either to lose male customers' business or risk license revocation---she suffered Article III injury-in-fact and was the ``obvious claimant'' to vindicate the concomitant rights of 18--20-year-old men whose market access would be materially impaired.

On the merits (Parts II-A through II-C), Brennan synthesized \textit{Reed v. Reed}, \textit{Frontiero v. Richardson}, \textit{Weinberger v. Wiesenfeld}, and \textit{Stanton v. Stanton} into a crisp, two-part test: gender classifications ``must serve important governmental objectives and must be substantially related to achievement of those objectives'' (\href{https://supreme.justia.com/cases/federal/us/429/190/}{429 U.S. at 197}). Accepting for argument's sake that Oklahoma's objective was traffic safety---an undeniably important state interest---Brennan dismantled the State's empirical showing. The most focused data, 1973 Oklahoma arrest statistics, showed 427 male versus 24 female DUI arrests and 966 male versus 102 female drunkenness arrests among 18--20-year-olds, translating to roughly 2\% of young men and 0.18\% of young women arrested for alcohol-related driving offenses. Even if statistically non-trivial, a 2\% correlation was an ``unduly tenuous `fit''' to justify using maleness as a legislative proxy. Brennan further catalogued severe methodological flaws: the surveys did not isolate ``nonintoxicating'' 3.2\% beer from stronger liquor; arrest data could be skewed by police chivalry toward young women; an Oklahoma City roadside survey showed only modest gender gaps in recent drinking (16.5\% of males vs. 11.4\% of females) and blood-alcohol levels above 0.01\% (14.6\% vs. 11.5\%); and the statute prohibited only \textit{sale} to young men, leaving them free to drink 3.2\% beer bought by female peers.

Finally, in Part II-D, Brennan traced the Twenty-first Amendment back to the Wilson and Webb-Kenyon Acts, concluding that \S\ 2 was intended primarily to create an exception to the Commerce Clause, not to qualify individual rights under the Bill of Rights or the Fourteenth Amendment. Relying on \textit{Wisconsin v. Constantineau} and \textit{Moose Lodge No. 107 v. Irvis}, Brennan held that the Twenty-first Amendment cannot save invidious gender discrimination, warning in footnote 22 that deferring to aggregate demographic drinking statistics under the Twenty-first Amendment would logically permit States to discriminate in alcohol sales by race, religion, or ethnicity.

\subsection{Concurring Opinion (Justice Powell)}

Justice Powell joined Brennan's opinion but wrote separately to express discomfort with formalizing a third, ``middle-tier'' level of equal protection scrutiny (\href{https://supreme.justia.com/cases/federal/us/429/190/}{429 U.S. at 210--211}). Acknowledging that ``candor compels the recognition that the relatively deferential `rational basis' standard of review normally applied takes on a sharper focus when we address a gender-based classification,'' Powell preferred framing the inquiry under \textit{Reed}'s ``fair and substantial relation'' test rather than proliferating formal tiers. Even crediting Oklahoma's statistics as showing that young men drove and drank more than young women, Powell viewed the case as ``relatively easy'' because a three-year gender gap that was ``so easily circumvented as to be virtually meaningless'' failed to bear a fair and substantial relation to highway safety.

\subsection{Concurring Opinion (Justice Stevens)}

Justice Stevens joined the Court's opinion while rejecting tiered scrutiny altogether: ``There is only one Equal Protection Clause. It requires every State to govern impartially'' (\href{https://supreme.justia.com/cases/federal/us/429/190/}{429 U.S. at 211--214}). Rather than sorting cases into mechanical tiers, Stevens asked whether the State's justification overcame the offensiveness of the classification. Here, Oklahoma's gender line was objectionable because it rested on an ``accident of birth,'' preserved a nearly extinct tradition of differential age-of-majority rules, and was physically ``perverse'' (since males, being generally heavier, can consume more alcohol before impairment). While the classification was not totally irrational, Stevens doubted traffic safety was the legislature's true purpose and found it unacceptable to impose a blanket disability on 100\% of young men---``visiting the sins of the 2\% on the 98\%''---for at most a minimal deterrent effect.

\subsection{Statement Concurring in Part (Justice Blackmun)}

Justice Blackmun issued a two-sentence statement joining all of the Court's opinion except Part II-D, while agreeing that the Twenty-first Amendment did not save the Oklahoma statute (\href{https://supreme.justia.com/cases/federal/us/429/190/}{429 U.S. at 214}). By withholding his join from Part II-D, Blackmun avoided endorsing the majority's sweeping curtailment of Twenty-first Amendment authority beyond what was necessary to decide the equal protection claim.

\subsection{Concurring in the Judgment (Justice Stewart)}

Justice Stewart agreed that Whitener had third-party standing under \textit{Eisenstadt} and \textit{Griswold}, and concurred in the judgment without adopting the majority's intermediate scrutiny formula (\href{https://supreme.justia.com/cases/federal/us/429/190/}{429 U.S. at 214--215}). Emphasizing that States possess broad power over alcohol under the Twenty-first Amendment, Stewart nonetheless reiterated from his \textit{California v. LaRue} concurrence that the Amendment does not empower a State to act with ``total irrationality or invidious discrimination.'' Because Oklahoma's statistics wholly failed to show that 3.2\% beer was more dangerous in the hands of an 18--20-year-old man than a woman of the same age, the law was totally irrational under \textit{Reed}.

\subsection{Dissenting Opinion (Chief Justice Burger)}

Chief Justice Burger dissented on both standing and the merits (\href{https://supreme.justia.com/cases/federal/us/429/190/}{429 U.S. at 215--217}). First, he rejected Whitener's \textit{jus tertii} standing as an unwarranted expansion of Article III exceptions: unlike in \textit{Barrows v. Jackson} or \textit{Griswold}, Oklahoma men aged 18--20 faced no barrier to suing in their own right (Craig's claim failed only because he turned 21), and equating a ``saloonkeeper'' with a physician-patient relationship ``borders on the ludicrous.'' On the merits, citing \textit{San Antonio Independent School District v. Rodriguez}, Burger argued that because gender is not a textually protected category or suspect class, only traditional rational-basis review under \textit{McGowan v. Maryland} applied. Since eight Justices recognized the law was not wholly irrational, the Court had no warrant to strike it down merely because it seemed unwise or ``a bit foolish.''

\subsection{Dissenting Opinion (Justice Rehnquist)}

Justice Rehnquist delivered a sharp doctrinal and statistical dissent (\href{https://supreme.justia.com/cases/federal/us/429/190/}{429 U.S. at 217--228}). First, he argued that heightened scrutiny was wholly unjustified where a statute disadvantaged \textit{men}, who lacked any history of pervasive discrimination or political powerlessness, and where the burdened interest (buying 3.2\% beer) was trivial. Second, he attacked Brennan's intermediate scrutiny test as conjured ``out of thin air'' without textual or precedential grounding, warning that terms like ``important'' and ``substantially'' were ``so diaphanous and elastic'' as to invite subjective judicial policymaking.

Third, under traditional rational-basis review (\textit{McGowan} and \textit{Williamson v. Lee Optical}), Rehnquist argued that the majority subjected the Oklahoma Legislature to ``the judicial equivalent of a doctoral examination in statistics.'' Critically, Rehnquist pointed out a logical flaw in Brennan's 2\% argument: equal protection rationality does not ask whether a large percentage of the regulated group exhibits the targeted trait (an inquiry that would invalidate even a gender-neutral 21-year-old drinking age, since only about 1\% of all 18--20-year-olds were arrested for DUI), but rather whether the \textit{relative incidence} of the hazard is sufficiently higher in the included class (an 18-to-1 male-to-female DUI arrest ratio, or 2\% vs. 0.18\%) to justify differential treatment. Because young men drove more miles, preferred beer more often, and had far higher arrest and casualty rates, Oklahoma's law was plainly rational.

\section{Examples: Future Applications}

\begin{enumerate}
  \item \textbf{Same Side (Unconstitutional):} A state law sets the minimum age to obtain a commercial motorcycle license at 21 for men and 18 for women, citing state highway data showing that 18--20-year-old men account for 88\% of motorcycle fatalities, even though less than 1.5\% of young male riders are involved in serious crashes. \textit{Why:} Like Oklahoma's beer statute, this facial sex classification uses maleness as a crude, overbroad proxy for reckless driving, burdening the vast majority of law-abiding young men when gender-neutral measures (such as graduated licensing or mandatory safety courses) directly target driving risk.
  \item \textbf{Same Side (Unconstitutional):} A state statute requires male applicants for state-subsidized nursing scholarships to sign a five-year post-graduation service commitment while requiring only a two-year commitment from female applicants, based on aggregate labor statistics showing male nurses change careers at slightly higher rates. \textit{Why:} Although retaining nurses is an important governmental objective, imposing a heavier legal burden on all male applicants based on aggregate statistical generalizations fails intermediate scrutiny's demand for a substantial, non-stereotyped fit.
  \item \textbf{Opposite Side (Constitutional):} A state enacts a uniform law raising the minimum purchase and possession age for all alcoholic beverages---including 3.2\% beer---to 21 for both men and women. \textit{Why:} As Justice Brennan noted in footnote 24 of \textit{Craig}, a gender-neutral age cutoff does not classify by sex and triggers only deferential rational-basis review, which traffic-safety data easily satisfy (as later reflected nationwide under the National Minimum Drinking Age Act of 1984).
  \item \textbf{Opposite Side (Constitutional):} A state public-health statute requires hospitals to offer free, routine prenatal blood screening specifically to pregnant patients, without offering identical prenatal tests to biological males. \textit{Why:} Because the classification tracks an actual biological difference rather than an archaic stereotype or loose statistical proxy, it bears a direct, substantial relationship to the important governmental objective of maternal and infant health (consistent with the Court's later recognition in \textit{Nguyen v. INS} and \textit{United States v. Virginia} that physical differences between the sexes are enduring).
  \item \textbf{Fence-Sitter (Borderline):} A state insurance regulator expressly permits auto insurers to charge 18--24-year-old male drivers 25\% higher premiums than female drivers of the same age and driving record, relying on comprehensive actuarial data showing male drivers in that bracket generate nearly twice the collision-claim payouts per mile driven. \textit{Why:} On one hand, actuarial pricing directly ties a much stronger empirical correlation to financial risk-pooling rather than criminalizing conduct, and private insurers lack a perfect individual proxy for first-time drivers. On the other hand, if state action is established through statutory mandate or rate-setting, \textit{Craig} warns against relying on sex-based aggregate statistics to penalize the safe majority of young men when gender-neutral proxies (such as telematics, mileage, or clean-record discounts) are available.
\end{enumerate}

\section{Critique}

\textit{Craig v. Boren} is foundational to modern Equal Protection doctrine, yet it has drawn sustained criticism across the ideological spectrum.

\textbf{Logical and Statistical Weaknesses:} Justice Rehnquist's dissent exposed a genuine analytical vulnerability in Justice Brennan's statistical critique. By focusing on the low \textit{absolute} rate of male DUI arrests (2\%) rather than the \textit{relative} disparity between young men and young women (2\% vs. 0.18\%, an elevenfold rate difference), Brennan's reasoning proved too much: under that metric, age-based drinking laws would also fail because 98--99\% of 18--20-year-olds are never arrested for DUI. Scholars have noted that Brennan's real objection was normative rather than purely statistical---namely, that the Equal Protection Clause forbids using sex as a proxy even when a real statistical correlation exists, a principle that sits uneasily alongside a test nominally asking only whether the classification is ``substantially related'' to an important end. Additionally, the Court's decision to expand third-party vendor standing for a beer retailer raised enduring concerns about lax Article III gatekeeping when the directly affected class (18--20-year-old men) could easily bring a class action.

\textbf{Originalist and Textualist Perspective:} From an originalist and textualist standpoint---echoing Chief Justice Burger, Justice Rehnquist, and later scholars and jurists---the Fourteenth Amendment's Equal Protection Clause contains no textual mandate for a three-tiered scrutiny architecture, nor was it originally understood in 1868 to prohibit sex-based classifications (\href{https://scholar.smu.edu/cgi/viewcontent.cgi?article=4964&context=smulr}{SMU L. Rev.}; \href{https://www.heritage.org/the-constitution/report/originalism-and-conservatism-american-story}{Heritage Foundation}). Indeed, Section 2 of the Fourteenth Amendment expressly referenced ``male inhabitants'' in apportioning congressional representation, and the Nation later adopted the Nineteenth Amendment to secure women's suffrage. Originalists argue that inventing ``intermediate scrutiny'' with elastic terms like ``important'' and ``substantially related'' transferred policy judgments from democratic legislatures to unelected judges, effectively enacting much of the unratified Equal Rights Amendment through judicial fiat (\href{https://ojs.library.okstate.edu/osu/index.php/OKPolitics/article/view/930/832}{Oklahoma Politics, 55}).

\textbf{Progressive and Feminist Perspective:} Conversely, progressive and feminist constitutional scholars have criticized \textit{Craig} from the opposite direction---for stopping short of strict scrutiny and obscuring the anti-subordination purpose of sex-equality law. As Reva Siegel has argued, beneath the ``faux precision'' of intermediate scrutiny lies deep doctrinal confusion about whether sex discrimination is less grave than race discrimination or simply different in kind (\href{https://law.yale.edu/sites/default/files/documents/pdf/Faculty/Siegel_SheThePeople.pdf}{Siegel, \textit{She the People}, Yale L.J.}). Because \textit{Craig} settled on an intermediate tier in a low-stakes case involving young men buying 3.2\% beer, it institutionalized a symmetrical ``anti-classification'' rule that treats benign or remedial laws aiding women with the same skepticism as laws entrenching patriarchy, while still leaving women with formally weaker constitutional protection than strict scrutiny provides for race (a gap only partially narrowed twenty years later by \textit{United States v. Virginia}'s ``exceedingly persuasive justification'' standard).

\section{Key Quotations}

\begin{quote}
To withstand constitutional challenge, previous cases establish that classifications by gender must serve important governmental objectives and must be substantially related to achievement of those objectives.
\end{quote}
\noindent --- Justice William J. Brennan, Jr., Opinion of the Court (\href{https://supreme.justia.com/cases/federal/us/429/190/}{429 U.S. at 197}).

\begin{quote}
It is unrealistic to expect either members of the judiciary or state officials to be well versed in the rigors of experimental or statistical technique. But this merely illustrates that proving broad sociological propositions by statistics is a dubious business, and one that inevitably is in tension with the normative philosophy that underlies the Equal Protection Clause.
\end{quote}
\noindent --- Justice William J. Brennan, Jr., Opinion of the Court (\href{https://supreme.justia.com/cases/federal/us/429/190/}{429 U.S. at 204}).

\begin{quote}
There is only one Equal Protection Clause. It requires every State to govern impartially. It does not direct the courts to apply one standard of review in some cases and a different standard in other cases.
\end{quote}
\noindent --- Justice John Paul Stevens, Concurring (\href{https://supreme.justia.com/cases/federal/us/429/190/}{429 U.S. at 211--212}).

\begin{quote}
But even assuming some such slight benefit, it does not seem to me that an insult to all of the young men of the State can be justified by visiting the sins of the 2\% on the 98\%.
\end{quote}
\noindent --- Justice John Paul Stevens, Concurring (\href{https://supreme.justia.com/cases/federal/us/429/190/}{429 U.S. at 214}).

\begin{quote}
The Court's conclusion that a law which treats males less favorably than females ``must serve important governmental objectives and must be substantially related to achievement of those objectives'' apparently comes out of thin air. The Equal Protection Clause contains no such language, and none of our previous cases adopt that standard.
\end{quote}
\noindent --- Justice William H. Rehnquist, Dissenting (\href{https://supreme.justia.com/cases/federal/us/429/190/}{429 U.S. at 220}).


\end{document}
